\documentclass[11pt]{article}
\usepackage[margin=1in]{geometry}
\usepackage{enumitem}
% =============================================================================
% Preambles/header.tex — shared LaTeX/Beamer preamble for this project
%
% Use from a Beamer lecture by prepending:
%     \input{header}
% and compiling with TEXINPUTS=../Preambles:$TEXINPUTS (the /compile-latex
% skill does this automatically).
%
% PALETTE CONTRACT. Color names defined here MUST match the names in
% Quarto/theme-template.scss so Beamer and Quarto renderings use the same
% palette. The scripts/check-palette-sync.sh script enforces this.
%
% To customize: edit the HEX values below AND the matching $variable in
% Quarto/theme-template.scss, then run scripts/check-palette-sync.sh.
% =============================================================================

% -----------------------------------------------------------------------------
% Engine detection + encoding
%
% Our /compile-latex skill uses XeLaTeX, where inputenc is unnecessary and
% can produce encoding warnings. Only load inputenc under pdfTeX.
% -----------------------------------------------------------------------------
\usepackage{iftex}
\ifPDFTeX
  \usepackage[utf8]{inputenc}
\fi

\usepackage{amsmath, amssymb, amsthm}
\usepackage{booktabs}
\usepackage{graphicx}

% -----------------------------------------------------------------------------
% PALETTE — mirrors Quarto/theme-template.scss variable names.
% Load xcolor and define colors BEFORE anything that references them
% (notably hyperref, hypersetup, and the Beamer color assignments).
% -----------------------------------------------------------------------------
\usepackage{xcolor}

\definecolor{primary-blue}{HTML}{012169}      % headings, accents
\definecolor{primary-gold}{HTML}{B9975B}      % emphasis, borders
\definecolor{highlight-yellow}{HTML}{F2A900}  % markers, alerts
\definecolor{light-bg}{HTML}{E8EDF5}          % title / section backgrounds
\definecolor{jet}{HTML}{1A1A1A}               % body text

% Semantic colors (used in diagrams and annotations)
\definecolor{positive}{HTML}{15803D}          % good / observed / identified
\definecolor{negative}{HTML}{B91C1C}          % bad / problematic / bias
\definecolor{neutral}{HTML}{525252}           % reference / context

% Highlight accents (match .hi-* classes in the SCSS)
\definecolor{hi-slate}{HTML}{314F4F}
\definecolor{hi-green}{HTML}{2E8B57}
\definecolor{hi-red}{HTML}{C41E3A}

% Semantic aliases so skill templates and snippets can write
% \textcolor{accent}{...} without hard-coding "primary-blue".
\colorlet{accent}{primary-blue}
\colorlet{accent2}{primary-gold}

% -----------------------------------------------------------------------------
% Hyperref — loaded AFTER xcolor + palette so link colors resolve.
% -----------------------------------------------------------------------------
\usepackage{hyperref}
\hypersetup{colorlinks=true, linkcolor=primary-blue, urlcolor=primary-blue,
            citecolor=primary-blue, pdfborder={0 0 0}}

% -----------------------------------------------------------------------------
% Course code — set once per deck (after \input{header}, before \begin{document})
% via \coursecode{CS401}. Shown in the footer on every slide so a deck's course
% is identifiable without opening the title page. Empty by default.
% -----------------------------------------------------------------------------
\makeatletter
\newcommand{\coursecode}[1]{\gdef\@coursecodetext{#1}}
\gdef\@coursecodetext{}
\makeatother

% -----------------------------------------------------------------------------
% Beamer theme assignments (only apply under Beamer).
%
% We detect Beamer via \@ifclassloaded{beamer}, which is the documented,
% non-internal way. This must be wrapped in \makeatletter/\makeatother
% because \@ifclassloaded contains an @-catcode character.
% -----------------------------------------------------------------------------
\makeatletter
\@ifclassloaded{beamer}{%
  \usefonttheme{professionalfonts}
  \setbeamercolor{structure}{fg=primary-blue}
  \setbeamercolor{frametitle}{fg=primary-blue}
  \setbeamercolor{title}{fg=primary-blue}
  \setbeamercolor{subtitle}{fg=primary-gold}
  \setbeamercolor{author}{fg=primary-blue}
  \setbeamercolor{institute}{fg=neutral}
  \setbeamercolor{date}{fg=primary-gold}
  \setbeamercolor{itemize item}{fg=primary-blue}
  \setbeamercolor{itemize subitem}{fg=primary-gold}
  \setbeamercolor{alerted text}{fg=primary-gold}
  \setbeamercolor{block title}{fg=white, bg=primary-blue}
  \setbeamercolor{block body}{bg=light-bg}
  % Semantic block variants: block=definition/concept (blue), exampleblock=
  % worked example (green/positive), alertblock=key takeaway or caution (gold).
  % Keep to <=2 colored boxes per slide across ALL three variants (INV-7).
  \setbeamercolor{block title example}{fg=white, bg=positive}
  \setbeamercolor{block body example}{bg=light-bg}
  \setbeamercolor{block title alerted}{fg=white, bg=primary-gold}
  \setbeamercolor{block body alerted}{bg=light-bg}

  % Minimal footer: course code (left) + page X / N (right), no navigation clutter
  \setbeamertemplate{navigation symbols}{}
  \setbeamertemplate{footline}{%
    \color{neutral}\footnotesize\hspace{0.7em}\@coursecodetext\hfill
    \insertframenumber/\inserttotalframenumber\hspace{0.7em}\vspace{0.3em}}
}{}
\makeatother

% -----------------------------------------------------------------------------
% TikZ setup — libraries the snippet gallery uses
% -----------------------------------------------------------------------------
\usepackage{tikz}
\usetikzlibrary{arrows.meta, positioning, calc, decorations.pathreplacing,
                fit, shapes.geometric, backgrounds}

% Shared TikZ styles — snippets and hand-written diagrams can pick these up
% via `\begin{tikzpicture}[every node/.style={...}]` or by naming specific
% styles (e.g., `\node[dag-node] (X) at (0,0) {X};`).
\tikzset{
  % Node styles for causal DAGs and similar
  dag-node/.style = {
    draw=primary-blue, fill=light-bg, rounded corners=2pt,
    minimum width=1.2cm, minimum height=0.9cm, align=center, font=\small
  },
  decision-node/.style = {
    draw=primary-gold, fill=white, diamond, aspect=1.8,
    minimum width=1.8cm, minimum height=1.2cm, align=center, font=\small,
    inner sep=1pt
  },
  % Edge styles — observed vs counterfactual
  observed-edge/.style = {-{Latex[length=2.5mm]}, thick, draw=jet},
  counterfactual-edge/.style = {-{Latex[length=2.5mm]}, thick, draw=neutral, dashed},
  confound-edge/.style = {-{Latex[length=2.5mm]}, thick, draw=negative, dashed},
  % Data-point styles for plots
  observed-dot/.style = {fill=primary-blue, draw=primary-blue, circle, inner sep=1.2pt},
  counterfactual-dot/.style = {fill=white, draw=neutral, circle, inner sep=1.2pt}
}

% -----------------------------------------------------------------------------
% Convenience macros
% -----------------------------------------------------------------------------
\newcommand{\muted}[1]{\textcolor{neutral}{#1}}
\newcommand{\key}[1]{\textcolor{primary-gold}{\textbf{#1}}}
\newcommand{\good}[1]{\textcolor{positive}{#1}}
\newcommand{\bad}[1]{\textcolor{negative}{#1}}

% Standout frame macro: full-bleed dark background for section transitions.
% Beamer-only (uses \begin{frame}/\setbeamercolor/\usebeamerfont) -- do not
% call from an article-class file (e.g. Notes/<CODE>/*-notes.tex). \muted,
% \key, \good, \bad, and \coursecode above are plain xcolor/gdef and are
% safe to use from either class; verified when Notes/article-class support
% was added.
% Expand: \transitionslide{Your transition title}
\newcommand{\transitionslide}[1]{%
  \begingroup
  \setbeamercolor{background canvas}{bg=primary-blue}
  \begin{frame}[plain]
    \centering
    \vfill
    {\usebeamerfont{title}\color{white}\Large #1\par}
    \vfill
  \end{frame}
  \endgroup
}

% Section divider frame (Beamer-only), an alternative to \transitionslide with a
% darker two-line style: near-black (jet) background, a small white label line
% above a larger gold title, and a thin gold rule along the bottom edge.
% Expand: \sectiondivider{Section 2}{Pointers}
\newcommand{\sectiondivider}[2]{%
  \begingroup
  \setbeamercolor{background canvas}{bg=jet}
  \begin{frame}[plain]
    \vspace*{3.0cm}
    {\color{white}\ttfamily\large #1\par}
    \vspace{0.5cm}
    {\color{primary-gold}\LARGE #2\par}
    \begin{tikzpicture}[remember picture, overlay]
      \draw[primary-gold, line width=2.5pt]
        ($(current page.south west)+(0,0.1)$) -- ($(current page.south east)+(0,0.1)$);
    \end{tikzpicture}
  \end{frame}
  \endgroup
}

% -----------------------------------------------------------------------------
% Channel watermark — "Concepts That Click" (YouTube).
%
% Article-class files (CompetitiveExam Q/A sets, Notes, Assignments) get a
% light diagonal draft watermark on every page plus a centered footer naming
% the channel. Beamer decks get a subtle TikZ background watermark instead
% (the footline already carries the course code). Centralized here so every
% MCQ PDF is branded without per-file edits.
% -----------------------------------------------------------------------------
\makeatletter
\@ifclassloaded{article}{%
  \usepackage{draftwatermark}
  \SetWatermarkText{Concepts That Click}
  \SetWatermarkScale{1}
  \SetWatermarkFontSize{2cm}
  \SetWatermarkLightness{0.86}
  \SetWatermarkAngle{45}
  \usepackage{fancyhdr}
  \pagestyle{fancy}
  \fancyhf{}
  \fancyfoot[C]{\footnotesize\textcolor{neutral}{Concepts That Click~|~YouTube: \href{https://www.youtube.com/@concepts.thatclick}{@concepts.thatclick}}}
  \renewcommand{\headrulewidth}{0pt}
  \renewcommand{\footrulewidth}{0pt}
  \fancypagestyle{plain}{%
    \fancyhf{}%
    \fancyfoot[C]{\footnotesize\textcolor{neutral}{Concepts That Click~|~YouTube: \href{https://www.youtube.com/@concepts.thatclick}{@concepts.thatclick}}}%
    \renewcommand{\headrulewidth}{0pt}%
    \renewcommand{\footrulewidth}{0pt}%
  }
}{}
\@ifclassloaded{beamer}{%
  \setbeamertemplate{background}{%
    \begin{tikzpicture}[remember picture, overlay]
      \node[rotate=45, opacity=0.055, align=center]
        at (current page.center)
        {\Huge\textcolor{neutral}{Concepts That Click}};
    \end{tikzpicture}%
  }
}{}
\makeatother 
\coursecode{JKSSB}
\renewcommand{\thesection}{37.\arabic{section}}
\newtheorem{example}{Example}[section]
\newenvironment{definitionbox}[1]{%
  \par\noindent\textbf{Definition (#1).}\ \itshape
}{\par\vspace{0.3em}}
\title{Email Basics --- Detailed Lecture Notes}
\author{JKSSB Computer Awareness}
\date{\today}

\begin{document}
\maketitle

\section{What email is}
Electronic mail (email) is a message sent from one electronic mailbox to
another over a network such as the Internet. Each user owns a mailbox
identified by a unique email address. The sender composes a message, directs
it with the recipient fields, and sends it; the message then waits in each
recipient's mailbox until it is opened. Because delivery does not require both
parties to be present at the same time, email is called a store-and-forward
service.

\begin{definitionbox}{Email address}
A unique identifier of a mailbox, written as a user name, then the @ symbol,
then the domain name, for example \texttt{assistant@gmail.com}.
\end{definitionbox}

\begin{center}
\begin{tabular}{p{0.24\textwidth}p{0.66\textwidth}}
\toprule
\textbf{Term} & \textbf{Meaning in this lesson} \\
\midrule
User name & The part of the address left of the @; the mailbox on the server. \\
Domain name & The part of the address right of the @; the mail server holding the mailbox. \\
To / Cc / Bcc & Recipient fields: primary recipient, visible copy, hidden copy. \\
Subject / body & The topic line and the main written content. \\
Attachment & A file sent along with the message. \\
Inbox / Sent / Drafts & Folders for received, sent, and unfinished mail. \\
Spam / Trash & Folders for filtered junk and for deleted messages. \\
Webmail / email client & Mail used in a browser vs mail used in an installed app. \\
\bottomrule
\end{tabular}
\end{center}

The rest of this lesson uses these terms in one consistent sense. The
\textbf{user name} is whatever stands left of the @ and the \textbf{domain
name} is whatever stands right of it. \textbf{To} is the main recipient,
\textbf{Cc} (Carbon Copy) is a visible additional recipient, and \textbf{Bcc}
(Blind Carbon Copy) is a hidden additional recipient. \textbf{Subject} is the
topic line, \textbf{body} is the message text, and an \textbf{attachment} is a
separate file carried with the message.

\section{The email address}
Every email address contains exactly one @ symbol, read as ``at''. The @
separates the user name on the left from the domain name on the right. In
\texttt{assistant@gmail.com}, \texttt{assistant} is the user name and
\texttt{gmail.com} is the domain name. The domain name identifies which mail
server holds the mailbox; the user name identifies which mailbox on that
server. This is the point tested when an item asks what the @ separates: the
answer is the user name and the domain name (PYQ-35).

\begin{definitionbox}{Domain name}
The part of an email address right of the @ that identifies the mail server or
organisation holding the mailbox, such as \texttt{gmail.com}.
\end{definitionbox}

Addresses are written without spaces, and the @ appears exactly once.
Anything before it belongs to the user name; anything after it belongs to the
domain. Options that say the @ separates the domain and the host, or two
parts of the domain, are therefore wrong.

\section{Recipient fields: To, Cc and Bcc}
The three recipient fields decide who receives the message and who can see
whom. \textbf{To} names the primary recipient, the person the message is
mainly for. \textbf{Cc} (Carbon Copy) names additional recipients who receive
a visible copy; every recipient can see the To and Cc lists. \textbf{Bcc}
(Blind Carbon Copy) names additional recipients whose addresses are hidden
from the other recipients.

\begin{definitionbox}{Bcc}
A recipient field whose addresses are hidden from the other recipients of the
same message; each Bcc recipient is concealed from the To, Cc, and fellow Bcc
recipients.
\end{definitionbox}

The word ``blind'' refers to this concealment. A Bcc recipient still receives
the message normally. The other recipients simply cannot see that address.
The sender, who composed the message, always knows the full list. To and Cc
addresses, by contrast, are visible to everyone who receives the message.
This is the distinction behind the common trap: Bcc hides recipients from each
other, not from the sender, and it does not hide the To or Cc lists from
anyone (TRAP-16).

\begin{example}
The officer sends a notice To the assistant and Cc the clerk, with the head
office on Bcc. The assistant and the clerk see each other's addresses but
cannot see the head office address. The head office receives the same notice.
The sender (the officer) knows all three addresses.
\end{example}

For the examination, read every Bcc option word by word. An option that says
Bcc hides the sender is wrong; an option that says Bcc hides the To list from
the recipients is wrong; the correct statement is that Bcc recipients are
hidden from each other while To and Cc remain visible to all.

\section{Subject, body and attachments}
A message has two content parts and an optional file part. The
\textbf{subject} is a single short line stating the topic, such as ``Leave
application for 12 October''. It is displayed in the message list before the
message is opened, so a clear subject helps the recipient identify the
message. The \textbf{body} is the main written content: the greeting, the
message itself, and the closing.

An \textbf{attachment} is a file sent along with the message, usually shown
with a paperclip icon. Common attachments are documents, spreadsheets,
presentations, PDFs, and images. Attaching a file sends a copy to the
recipient; the sender keeps the original. Two practical points matter for the
examination. First, mail servers impose a size limit, so a very large file is
rejected or must be shared as a link instead. Second, an attachment can carry
a virus, so an unexpected attachment should not be opened blindly.

\section{Folders: Inbox, Sent, Drafts, Spam, Trash}
Received mail arrives in the \textbf{Inbox}. Every account has exactly one
Inbox; it is the default view when the mailbox is opened. A copy of every
message that is sent is kept in the \textbf{Sent} folder (sometimes labelled
Sent Items or Sent Mail). A message that has been started but not yet sent can
be saved in \textbf{Drafts}; it waits there until it is completed and sent.

\begin{definitionbox}{Drafts}
The folder holding composed messages that have been saved but not yet sent.
\end{definitionbox}

Unwanted or bulk mail is separated into the \textbf{Spam} folder (also called
Junk). The mail service filters suspected junk out of the Inbox automatically;
a wanted message found there should be moved back or marked ``not spam''.
Messages the user deletes go to the \textbf{Trash} folder (also called Deleted
Items). They remain recoverable until the folder is emptied; emptying Trash
or Spam permanently removes those messages.

The five folders are most easily confused in pairs. Inbox receives while Sent
keeps sent copies; Drafts holds work not yet sent while Sent holds work
already sent; Spam holds filtered junk while Trash holds what the user deleted.
No other folder does another's job: drafts never send themselves, sent copies
never return to drafts, and spam is not the delete bin.

\section{Reply, Reply All and Forward}
Three actions dispose of a received message. \textbf{Reply} composes an answer
addressed only to the sender. \textbf{Reply All} composes an answer addressed
to the sender plus all the To and Cc recipients of the original message.
\textbf{Forward} sends the received message onward to new recipients chosen
by the forwarder, usually with a short added note on top.

\begin{definitionbox}{Reply All}
A response addressed to the sender and to every visible (To and Cc)
recipient of the original message; Bcc recipients of the original are never
re-added.
\end{definitionbox}

Two consequences follow. First, Reply All can only reach the addresses visible
on the original message, so hidden Bcc addresses are never included. Second,
a forwarded message carries the original content (and often its attachments)
inside the new message; the forwarder becomes the new sender. For the
examination, match the verb to the action: ``answers only the sender'' is
Reply; ``answers everyone visible'' is Reply All; ``sends it onward to new
recipients'' is Forward.

\section{Webmail versus an email client}
Email can be used in two ways. \textbf{Webmail} means using mail through a
web browser, with the mailbox kept on the provider's server; Gmail, Yahoo
Mail, and Outlook.com are webmail services. An \textbf{email client} is an
installed application used for mail, such as Microsoft Outlook or Mozilla
Thunderbird; it can keep or download copies on the local machine.

Closely tied to the client are the incoming-mail protocols. \textbf{POP3}
(Post Office Protocol version 3) downloads mail to the local machine, usually
removing it from the server. \textbf{IMAP} (Internet Message Access Protocol)
keeps mail synced on the server, so a phone, an office computer, and a laptop
all show the same folders. Outgoing mail in both arrangements is normally
sent with \textbf{SMTP} (Simple Mail Transfer Protocol). The one-line
contrast to remember is: POP3 downloads, IMAP syncs, SMTP sends.

\begin{example}
The assistant reads office mail in two places: through Gmail in the browser
(webmail) and through the Outlook app on the office computer (email client)
configured with IMAP, so a message read on one device shows as read on the
other. Had the client used POP3, the downloaded message would sit on that one
computer instead of staying synced on the server.
\end{example}

\section{Exam retrieval and common confusions}
Six distinctions prevent most errors on this topic.
\begin{enumerate}
  \item The @ separates the \textbf{user name} (left) from the \textbf{domain
    name} (right). It does not separate domain and host (PYQ-35, TRAP-28).
  \item \textbf{Bcc} hides recipients from \textbf{each other}; the sender
    always knows, and To/Cc stay visible to all (TRAP-16).
  \item \textbf{Cc} is Carbon Copy (visible); \textbf{Bcc} is Blind Carbon Copy
    (hidden). Do not swap ``blind'' and ``visible''.
  \item \textbf{Reply} answers the sender; \textbf{Reply All} answers all
    visible recipients; \textbf{Forward} sends the message onward.
  \item \textbf{Drafts} holds unsent work, \textbf{Sent} holds sent copies,
    \textbf{Inbox} holds received mail, \textbf{Spam} holds filtered junk,
    and \textbf{Trash} holds deleted mail.
  \item \textbf{Webmail} runs in a browser; an \textbf{email client} is
    installed software. \textbf{POP3} downloads mail to the machine while
    \textbf{IMAP} keeps it synced on the server.
\end{enumerate}

\subsection*{Exercises}
\begin{enumerate}[label=\arabic*.]
  \item In the address \texttt{assistant@gmail.com}, identify the user name
    and the domain name. What does the @ separate?
  \item Distinguish To, Cc, and Bcc by stating who can see whom.
  \item Distinguish the subject, the body, and an attachment.
  \item State which folder holds each of the following: an unfinished message,
    a copy of a sent message, a newly arrived message, filtered junk, a
    deleted message.
  \item Distinguish Reply, Reply All, and Forward.
  \item Distinguish webmail from an email client, and POP3 from IMAP.
\end{enumerate}

\subsection*{Solutions}
\begingroup\small
\begin{enumerate}[label=\arabic*.]
  \item The user name is \texttt{assistant} and the domain name is
    \texttt{gmail.com}. The @ separates the user name (left) from the domain
    name (right).
  \item To is the primary recipient and Cc is a visible copy; both lists are
    seen by all recipients. Bcc is a hidden copy: each Bcc address is concealed
    from the other recipients, while the sender knows the full list.
  \item The subject is the short topic line shown in the message list; the
    body is the main written content; an attachment is a separate file carried
    with the message (shown with a paperclip icon).
  \item Unfinished: Drafts. Sent copy: Sent. Newly arrived: Inbox. Filtered
    junk: Spam (junk). Deleted: Trash (deleted items), recoverable until
    emptied.
  \item Reply answers only the sender; Reply All answers the sender plus all
    visible To and Cc recipients (never Bcc); Forward sends the received
    message onward to new recipients.
  \item Webmail is mail used through a browser with the mailbox on the
    server (Gmail, Yahoo Mail, Outlook.com); an email client is installed
    software (Outlook, Thunderbird). POP3 downloads mail to the local
    machine, usually removing it from the server; IMAP keeps mail synced on
    the server so all devices show the same folders.
\end{enumerate}
\endgroup
\end{document}
